\chapter{Як машина вчиться}
% full version: chapters/06_dofa.tex

У всіх досі розглянутих схемах силу зв'язків добирав інженер. У живому мозку інженера немає. Зв'язки мають налаштовуватися самі, під час роботи, --- і так, щоб машина робила щось корисне. Це і є навчання.

Головне обмеження звучить дуже суворо. Кожне з'єднання змінює себе саме, і знати воно може лише те, що відбувається поруч: чи працює відправник, чи працює отримувач і які речовини до нього доходять. Ніхто не може надіслати з'єднанню особисту поправку. Це виключає головний метод навчання штучних нейромереж, де похибка пробігає мережу у зворотний бік і кожен зв'язок отримує свою власну поправку.

\section*{Разом --- отже, пов'язані}

Найпростіше правило: якщо відправник і отримувач працюють разом, зв'язок між ними міцнішає. Воно локальне й не потребує годинника. Таке правило, якщо додати до нього трохи забування, знаходить у потоці сигналів найголовніше --- те, що змінюється найсильніше. Є й версія, яка зважає на час клацань: зв'язок міцнішає, якщо відправник клацнув \emph{перед} отримувачем, тобто міг бути причиною, і слабшає, якщо після.

Але ці правила мають одну межу: вони вчать того, що \emph{часто}, а не того, що \emph{корисно}. З'єднання, яке бачить лише своїх сусідів, не знає, чи принесла дія сік чи біль.

\section*{Третій сигнал}

Його приносить дофамін --- радіостанція «краще, ніж очікувалося». Нехай зв'язок змінюється, лише коли збіглися три речі: працював відправник, працював отримувач і прийшов дофамін.

\pencilfig{6}{\pencilart{6}}{Зв'язок міцнішає, коли обидва нейрони працювали разом і прийшов третій сигнал: «вийшло краще, ніж чекали».}

Але ось складність: винагорода приходить не одразу, а за секунду-другу. На цей момент відправник і отримувач давно замовкли. З'єднанню потрібна пам'ять про те, що воно брало участь. І вона є: спільна робота залишає в з'єднанні мітку, як стікер, що сам відклеюється за секунду-дві. Прийшов дофамін, поки стікер тримається, --- зв'язок змінюється. Не прийшов --- мітка зникає безслідно. Це розв'язує задачу адресації: дофамін один на всіх, але змінюються лише ті з'єднання, на яких висить стікер.

\pencilfig{106}{%
\foreach \y/\d/\t in {0/0.9/{вчасно: зв'язок міцнішає},-1.3/3.3/{запізнився: нічого}}{
  \draw[pen] (0,\y) -- (4.6,\y);
  \foreach \s in {0.25,0.33}{\draw[pcBlue, line width=0.8pt] (\s,\y) -- (\s,\y+0.3);}
  \fill[pcAmber, opacity=0.25] (0.35,\y) -- plot[domain=0.35:4.5, samples=30] (\x,{\y+0.8*exp(-(\x-0.35)/0.8)}) -- (4.5,\y) -- cycle;
  \draw[penlite=pcAmber] plot[domain=0.35:4.5, samples=30] (\x,{\y+0.8*exp(-(\x-0.35)/0.8)});
  \draw[pcAmber!80!black, line width=0.9pt] (\d-0.1,\y+0.95) -- (\d+0.07,\y+0.62) -- (\d-0.07,\y+0.6) -- (\d+0.1,\y+0.22);
  \draw[arr=pcAmber!80!black] (\d+0.09,\y+0.27) -- (\d+0.11,\y+0.12);
  \node[lbl, anchor=west] at (4.7,\y+0.3) {\t};}
\node[lbl, text=pcBlue, anchor=south west] at (0.1,0.85) {працювали разом};
\node[lbl, text=pcAmber!80!black, anchor=south] at (2.3,0.35) {мітка тане};
\foreach \x/\l in {0.35/0,1.95/1,3.55/2}{\draw[pcGray, line width=0.5pt] (\x,-1.3) -- (\x,-1.38); \node[lbl, anchor=north] at (\x,-1.38) {\l~с};}
}{Спільна робота залишає в з'єднанні мітку, яка тане за секунду-дві. Дофамін змінює лише з'єднання з міткою.}

\section*{Шум як пошук}

Чому це правило веде до кращого, а не куди завгодно? Уявіть, що ви із зав'язаними очима шукаєте вершину пагорба. Ви робите випадковий крок, і хтось гукає «вище!» або «нижче!». Якщо «вище» --- закріплюєте крок. Крок за кроком ви піднімаєтеся, жодного разу не глянувши на мапу. Мозок робить те саме: нейрони трохи тремтять випадковим чином, дофамін повідомляє, чи стало краще, і зв'язки, що брали участь, зсуваються у вигідний бік. Шум, який заважав передавати повідомлення, тут стає пошуком.

Тепер зрозуміло, чому дофамін повідомляє не про винагороду, а про \emph{сюрприз}. Якби він просто кричав «сік!» після кожного успіху, він заразом закріплював би все, що робила мережа, --- і правильне, і неправильне. У нашій моделі така мережа справді розганяє свої зв'язки в тисячі разів, а іноді назавжди застрягає на гіршому виборі. Віднімання очікування робить навчання чесним.

\section*{Ціна одного проводу}

Широкомовне навчання має свою ціну. Сигнал один на всіх, а в ньому змішано шум усіх нейронів, що впливають на результат. Що більша мережа, то довше кожному з'єднанню виокремлювати свою частку. Тому дофамін, схоже, навчає лише порівняно невеликі кола, що обирають дії, а величезна кора вчиться переважно за місцевими правилами.

\section*{Хто рахує сюрприз}

Лишається питання: як дофамінові нейрони обчислюють «краще, ніж очікувалося»? Потрібен критик --- група нейронів, яка весь час оцінює, скільки доброго ще попереду. Сюрприз --- це отримана винагорода плюс те, наскільки підскочило очікування. Засвітилася лампочка --- очікування підстрибнуло, сплеск. Прийшов обіцяний сік --- винагорода є, але очікування тієї ж миті вичерпано, сюрпризу немає. Сік не прийшов --- очікування обвалилося даремно, провал. Усі три випадки з першого розділу виходять самі собою. Що дофамін поводиться саме так, виміряно. Як саме мозок обчислює цей сюрприз --- поки гіпотеза.

\fullversion{6}
